All 54 real-data fits completed (nine participants, three paired split/training
seeds, two arms). The following values average seeds within participant and
participants equally. Accuracy values are percentages; differences are percentage
points. Log loss is unitless and lower is better.

\begin{center}\begin{tabular}{@{}lrrrr@{}}\toprule
Model & Train BA & Validation BA & Gap & Val. log loss \\ \midrule
Frozen factors & 67.44 & 45.71 & 21.73 & 1.593 \\
Supervised factors & 72.86 & 47.47 & 25.40 & 1.570 \\
\bottomrule\end{tabular}\end{center}

\begin{center}\begin{tabular}{@{}lrrr@{}}\toprule
Participant & Frozen BA & Supervised BA & Difference \\ \midrule
A01 & 39.74 & 45.70 & +5.95 \\
A02 & 38.10 & 41.80 & +3.71 \\
A03 & 59.71 & 63.42 & +3.71 \\
A04 & 42.36 & 41.93 & -0.43 \\
A05 & 40.38 & 41.03 & +0.64 \\
A06 & 42.42 & 42.58 & +0.15 \\
A07 & 41.85 & 43.04 & +1.19 \\
A08 & 48.53 & 49.08 & +0.55 \\
A09 & 58.32 & 58.62 & +0.31 \\
\bottomrule\end{tabular}\end{center}

\begin{center}\begin{tabular}{@{}lrrr@{}}\toprule
Availability & Frozen BA & Supervised BA & Difference \\ \midrule
Full 22 & 45.71 & 47.47 & +1.75 \\
Static 16 & 37.94 & 40.03 & +2.08 \\
Dynamic 16 & 37.61 & 40.47 & +2.86 \\
Static 6 & 35.42 & 35.38 & -0.04 \\
Dynamic 6 & 35.52 & 36.37 & +0.85 \\
\bottomrule\end{tabular}\end{center}

\textbf{Primary finding.} Updating the factors improved full-input balanced
accuracy by 1.75 percentage points, with an exploratory participant-bootstrap
95\% interval of $[0.56,3.18]$ points. Eight of nine participant averages improved;
A04 declined by 0.43 points. However, the prespecified two-point mean-gain
threshold was not met. The appropriate verdict is a small, fairly consistent
benefit from supervised factors, with a failed practical screening criterion.
It is not a null result, nor evidence of a strong classification architecture.

Split/training-seed mean gains were $+1.80$, $+2.58$ and $+0.87$ points for
seeds 0, 1 and 2. Seventeen of the 27 individual participant/seed pairs improved,
eight declined and two tied (absolute tolerance $10^{-12}$ in balanced accuracy).
Training balanced accuracy increased by 5.42 points while validation increased
by 1.75, widening the mean train--validation gap from 21.73 to 25.40 points.
This is consistent with additional fitting capacity exceeding the gain in
generalization; it does not uniquely identify the cause. Full-input log loss
changed by $-0.023$ with interval $[-0.243,0.182]$, providing no clear improvement
in probabilistic prediction. Both mean full-input log losses exceed the uniform
four-class predictor's $\log 4\simeq1.386$, despite above-chance balanced accuracy.
This indicates a weakness in confidence quality. All 54 fits stopped before
the 250-epoch limit (75--121 epochs); extending the epoch budget is not supported
by these histories.

\textbf{Electrode-loss findings.} At 16 retained channels, gains were
$+2.08$ points for static masks (interval $[0.81,3.10]$) and $+2.86$ for dynamic
masks ($[1.08,4.47]$), each positive in eight participants. At six channels,
static masks changed by $-0.04$ points ($[-1.28,1.45]$) and dynamic masks by
$+0.85$ ($[-0.89,2.72]$). Thus this intervention helps moderate-loss accuracy,
but does not establish a severe-loss benefit. Mean log loss at six channels
increased by 0.068 (static) and 0.027 (dynamic), with both intervals crossing
zero. All these condition-specific findings are exploratory.

\textbf{Research decision.} Keep differentiable factor learning as a modestly
supported component, but do not promote this model as the accuracy baseline.
The 47.47\% validation balanced accuracy remains below the earlier spatial
Transformer's 63.71\% single-seed screen and 64.87\% three-seed study; those are
historical context, not matched contrasts in this study. The result supports
classification-informed factor adaptation at these fixed spectral features,
ranks and head, but cannot show that task adaptation alone repairs the
information bottleneck. It also does not test learned raw temporal features.

A useful next bounded comparison would train both frozen and supervised arms
with the same mixed-availability schedule, retaining this full-training study
as a control. That would test whether exposure to missing channels improves
the six-channel outcome. For the separate full-input accuracy objective,
test a simpler or more strongly regularized classifier on the same cores before
increasing factor ranks or Transformer depth. The larger train--validation gap
motivates this test but does not prove that classifier capacity is the cause.
Each change
requires its own declared identity and controls; neither experiment was run here.

\textbf{Verification and provenance.} The independent audit replays all 54
selected checkpoints on training and validation data, reconstructs all 1,134
availability evaluations, recalibrates each task's training-only initial factors,
and independently checks aggregate means and bootstrap intervals. Selected
supervised factors changed in all 27 fits, while frozen factors were unchanged;
paired initial full states and calibration histories match exactly. The software
checks comprise 25 new model/runner tests and 17 tensor temporal regression tests.
The numerical review checks the direct system against an explicit observed-entry
design and verifies gradients. Historical source/configuration/workflow/result
preservation checks cover 1,122 files.

All 85 Python source hashes, 64 input hashes, configuration and package versions
are recorded in the manifest. The study identity is
\begin{center}{\scriptsize\ttfamily
479db2ec18e7f1e1b5d1e5255c5a200c3df7d59b7d8a73936b4415252af8ed57}
\end{center}
The machine-readable results and independent audit are under
\code{results/supervised-tucker-v1/report/}. Fits used one CPU thread each;
cohort execution used up to four independent processes. Elapsed times are
execution records, not controlled speed benchmarks. Checkpoint selection and
reporting used validation only, so an independent session or dataset remains
necessary before a confirmatory accuracy claim.
